\chapter{Resumen ejecutivo}
La campaña EXEC\_46 separó tres cuestiones que habían quedado mezcladas en el diagnóstico inicial. La no-linealidad observada de $T_0(x)$ se volvió a evaluar con una distinción explícita entre pendientes \emph{within} y \emph{between}, y con especificaciones lineal y cuadrática de la curva entre posiciones. Las tablas de Step~5 documentan una reducción del residuo al cambiar la especificación, sin convertir esa reducción en una identificación causal.

La selección del primer fotón muestra una población Cherenkov dominante en el primer disparo de los extremos cercanos. La comparación de la fracción entre la campaña constante y la campaña v2 está en las tablas T2; los valores usados son macros generadas desde los CSV registrados.

El modelo óptico constante fue sustituido, para EJ--200 y EJ--204, por tablas dispersivas basadas en Huggins et al.~\cite{Huggins2026}. EJ--230 permanece constante por ausencia de un análogo medido. La campaña v2 fue validada frente a la campaña antigua con EJ--230 como control.

El observable de velocidad efectiva depende de la convención temporal y de la forma de pulso asumida. La comparación con el valor experimental debe leerse junto con la ausencia de una parametrización MUSIC/SAMPIC en las fuentes disponibles. Quedan abiertas la parametrización de la cadena electrónica y la descomposición completa del término geométrico del tiempo efectivo.
